\documentclass[11pt]{article}
\usepackage[margin=1in]{geometry}
\usepackage{microtype}
\usepackage{parskip}

\title{A Theory Toward the Structure of Prolonged Existence\\[0.35em]
\large A Concluding Articulation}
\author{Albert Jan van Hoek}
\date{June 2026}

\begin{document}
\maketitle

\section*{Introduction}

This is a concluding statement of a body of work I have developed under the name \emph{Evolution by Emergence}. I set it down not as a summary of the parts but as an articulation of the single thing they have turned out to be about.

The work began with a narrow question: whether the intelligence now appearing in machines is a novel and magical thing, or an instance of something reality already does. Pursued far enough, the answer reframed the whole enterprise. What I have been describing is not, in the end, a theory of evolution. Evolution is the name for how one substrate---life---changes. What this is a theory of is more general, and sits a level of abstraction higher: the structural conditions under which \emph{any} configuration persists, on any substrate, when the world it lives in will not hold still. I will call it, plainly, a \emph{theory toward the structure of prolonged existence}. The renaming is not cosmetic. Naming the work by its object---prolonged existence---rather than by one of its instances---biological evolution---is what frees it to span substrates without becoming an analogy stretched from biology.

The first move is a demotion. Natural selection is usually treated as the engine of adaptation; I treat it as one species of a wider genus. Above it sits a more general constraint---adaptive persistence---which holds that where no fixed configuration suffices, only systems that can reconfigure themselves persist. Selection is biology's way of meeting that constraint; belief revision is cognition's; elections and bankruptcy are some of the institutional ways. None need share a mechanism. They share a function.

\section*{The internal logic}

A theory earns the right to predict and explain not from its breadth but from its internal logic, and it is the logic, not the scope, that does the work here.

Begin with a gradient---any difference the world has not yet equalised. Structures arise that exploit it. What we call efficiency is best understood not as a quantity that tends to increase but as the lengthening of a becoming along the same gradient: a raw gradient gives one burst and collapses to equilibrium, while a structure that exploits it more efficiently keeps it productive for longer. Greater efficiency therefore buys greater duration---which is why the structures still present to be observed are disproportionately the efficient ones. This is a selection effect, not a trend. Nothing climbs toward efficiency; the inefficient simply do not last to be counted.

Where a part can be rewarded locally while the whole loses globally, a second thing follows, which I call \emph{capture}: parts tend to corner the shared substrate and exhaust it---the predator that eats all its prey, the monopoly that owns a dead market. Wherever that local/global gap exists, capture is a permanent failure mode. It never stops arising, and most instances are terminal: the substrate collapses, and the system dies or falls back. The correction is reliable, but its route is not foreordained. Finitude guarantees that a corrective arrives---if not the gentle one a system installs deliberately, then the brutal one that collapse installs for free. The only real choice is which. And the corrective settles, reliably, on whoever cannot externalise the cost, because for them the substrate is not something to offload; it is the ground they stand on.

Occasionally a capture is not suffered but \emph{domesticated}. A higher-level unit forms that constrains the captor while preserving its function---what was external capture becomes internal function. These are the major transitions: the captured part surrenders some autonomy, and a new unit appears that occupies adaptive space its components could not reach alone. From this follows the account of complexity, and it is deliberately non-teleological. Complexity is not where the process is going; it is what remains when capture is solved instead of suffered. The successes are rare and reversible, and every new level is also a fresh arena for the same problem one floor up. There is no smooth ascent---only a graveyard of terminal captures with a few domesticated ones still standing, each of which became the ground the next could build on.

Read along the vertical, persistent structures form a \emph{commons}. Each exists because some prior possibility opened, and each, by existing, opens the next. The relation runs both ways, which is what makes it a commons and not a ladder: what a structure draws from the shared space it also returns to it, so that contributing to the commons---holding one's own gradient open efficiently---is the same act as extending what everything else can become. The value a persistent structure brings to the whole is, at bottom, time. Capture and commons are then two faces of one thing: capture is a part taking from the shared space faster than it returns; the commons is the shared space itself, the openness that capture threatens and that organised structure exists to keep open.

What differs across the commons is not whether a structure participates but its \emph{bandwidth}---how long it holds its solutions open, and how far it can recombine them. A crystal holds its solution only until conditions shift, then re-equilibrates and forgets; a lineage that inherits and recombines holds its solutions across conditions and stacks them. That difference is not a wall between the living and the non-living, or between selection and stability. It is a gradient within the shared space, and it is what sets the rate at which complexity accumulates.

Two further features keep the picture honest. The property that holds a commons together is best called \emph{integration} rather than organisation: interdependence that settles in through paths of competition and is held in place by the web it comes to depend on---integration without an integrator. Organisation, in the deliberate human sense, is only its special case at the high-deliberation end; emergence, in the weak sense of a resolution derivable only by running the process, is the general phenomenon. And every durable structure is layered: a slow, shielded core funds a fast, freely revised edge, and the shielding of the core is precisely what funds the fluidity of the edge. The strength of each corrective is a parameter set per layer, not once for the whole; where on the band of viable settings a system sits is a value choice to be revisited, not a number the theory can supply. None of the layers may be allowed to stop, for a frozen core that scaffolded one era becomes the constraint of the next.

\section*{What it adds}

Set beside the standard account of evolution, the addition is locatable rather than rhetorical. Where the modern synthesis foregrounds the replicator and treats the network as environment, this foregrounds the network and treats the replicator as one kind of participant. The web is moved from the edge to the centre.

It also clarifies what kind of unity is being claimed, and what kind is not. This is not one process. The upgrading of a gene and the upgrading of a legal text are incomparable as mechanisms; selection and deliberation share no machinery. What they share is a structural role---contributing duration to the commons, and being exposed to capture along the gradient. The unity is structural, not mechanical: many incomparable processes occupying one structure, wearing different masks.

\section*{What it lets us do}

Because the structure has an internal logic, it yields predictions and explanations rather than only descriptions.

Its cleanest current application is the question it began from. Artificial intelligence is a new substrate running the same structure: it is a new participant in the commons, and the prediction is one of form, not outcome. By the framework's own logic, a participant of unprecedented \emph{bandwidth} is precisely what can capture the shared space faster than the rest can mount a corrective. So the framework does not conclude that intelligence on a new substrate is safe because it belongs to the same process; it predicts the opposite---that such a participant is the sharpest capture risk the commons has met, which is the alignment problem stated in the framework's own terms. The answer it offers is its standing one: keep the commons open, lower the cost at which the rest can coordinate, and do not let the high-bandwidth part foreclose the shared space.

To the human-scale captures---the concentration of wealth in a few, the capture of a shared climate, durable asymmetries of role and power---the framework speaks diagnostically, not prophetically. It predicts the \emph{structure} of such problems reliably: that they are local/global gaps; that the corrective must come from those who bear the collapse; that capture, blocked in one channel, migrates to the next; that the cost is asymmetric. It does not predict the outcome, because its own logic is that most captures are terminal. For the climate in particular it predicts that the physical floor is the one sensor that cannot be jammed and the one that fires late, at full cost---a sobering forecast, not a reassuring one, and a more useful one for being honest.

\section*{What it does not give}

Three limits are part of the theory, not apologies for it.

It rests on a named value. That a system's persistence requires its corrigibility is logic; that we should build arrangements which keep the adaptive space open is a commitment, and the one premise the whole structure rests on and cannot itself supply. Without it the descriptions remain descriptions, with no \emph{ought} anywhere in them.

It offers no seal. Every corrective leaks, is gamed, is routed around; rebuilding them is the work, not evidence it was done wrong. A system that believed it had permanently solved capture would have committed the central error in its purest form---a configuration that escaped its own revision.

And it is no exception to itself. It is a configuration in the same adaptive space as everything it describes. If it is wrong it should be replaced; if it is partial, expanded; if it is useful, it has earned no more than provisional survival, for as long as it keeps passing the test. A theory immune to its own principle would be the exact form of capture the principle names.

What it claims, then, is modest in kind and not in reach: not that it explains everything, but that it names a structure with an internal logic that does real work---a frame that makes an aspect of reality already there somewhat clearer. Survival has never belonged to the strongest, nor to the most certain. It belongs to whatever keeps itself able to change, and, where there are many, to whatever can hold open the space in which the changing happens. Even this conclusion asks to be doubted. That is not a weakness of it. It is the argument.

\end{document}
